\chapter{Conclusions and Recommendations}
\label{ch:conclusion}

This final chapter draws together the work presented in the preceding chapters. It first
states the overall conclusion of the project and situates it against the objectives set out
in Chapter~\ref{ch:introduction}; it then itemises in detail what was actually delivered;
it reflects on the technical and teamwork lessons that emerged during development; it
discusses the system's honest limitations without overstating its maturity; and it finally
sets out a prioritised programme of future work and recommendations. Taken as a whole, the
chapter is intended to serve two audiences: an assessor seeking a candid appraisal of the
degree to which the project met its goals, and a future maintainer seeking a clear sense of
where the platform stands and where it could profitably be taken next.

\section{Conclusion}
\label{sec:conclusion-summary}

The objective of this graduation project was to design and build a single, coherent,
cloud-native platform that consolidates the fragmented operational concerns of a
multi-sport club into one secure and analytically rich system. Taking FC~Barcelona as
the working organisational model and \textit{Sportify} as the product brand, the
project set out to replace the patchwork of spreadsheets, e-mail threads and
disconnected SaaS tools described in Chapter~\ref{ch:introduction} with a unified platform
on which each professional role---from leadership through coaching, medical, scouting and
commercial staff down to players and fans---sees exactly the information relevant to its
responsibilities, and nothing more. The thesis underpinning the work was that the
structural problems of club management---data fragmentation, multi-discipline complexity
and analytical blindness---are not problems of any single feature but of architecture, and
that a properly decomposed, identity-centred, event-driven system could dissolve them at
their root.

This objective was met. The Multi-Sport Club Management System (MSCMS) is a working
platform built on a microservices architecture~\cite{microservices} of six domain
services and a three-part control plane---together seven independently deployable
Spring~Boot~3.5.6 / Java~21 services~\cite{springboot}---discovered through Eureka and
reached through a single Spring Cloud Gateway~\cite{springcloud}, fronted by a
Next.js~16 / React~19 application~\cite{nextjs,react}, and orchestrated as containers
through a single Docker~Compose stack~\cite{docker}. Every request is authenticated against
a Keycloak realm and verified as a JSON Web Token at the gateway~\cite{keycloak,jwt}, and
persistence follows the database-per-service pattern over PostgreSQL~\cite{postgres} with
asynchronous, event-driven coordination over Apache~Kafka~\cite{kafka}. The architectural
hypothesis was therefore not merely asserted but demonstrated in running software: each of
the structural problems identified at the outset is answered by a concrete architectural
mechanism---fragmentation by a shared identity and a single gateway origin, multi-discipline
complexity by sport-aware domain models, and analytical blindness by a reporting service
that aggregates across the previously siloed domains.

Beyond reproducing the operational baseline, the project advanced the system from a
predominantly CRUD-oriented tool into a genuinely interactive sports platform. The
headline outcomes are a database-backed competitions module that models La~Liga and the
new-format UEFA Champions League with a single league phase and a knockout bracket, fed by
real fixture data ingested from football-data.org~\cite{footballdata}; a broadcast-style
live-match engine that drives a match through its phases in real time, complete with
extra time and penalty shoot-outs for knockout ties; a clinical injury-management
workflow that follows each injured player through a staged recovery journey; a real-time
team messaging facility built over WebSocket~\cite{websocket}; an event-driven
notification centre; and consistent role-based access control~\cite{rbac} enforced at
identity, gateway and user-interface layers. Collectively, these capabilities
demonstrate that the architectural decisions taken early in the project---service
decomposition, a central identity provider, an asynchronous event backbone and a
REST-based contract between front-end and back-end~\cite{restapi}---scaled gracefully to
accommodate substantially richer behaviour without structural rework. It is worth
emphasising this last point: none of the later, behaviourally complex modules required a
new service, a new gateway route or a change to the authentication model. The competitions
module and the live-match engine were absorbed by the training-match service because they
share the match and line-up schema; the staged injury workflow was absorbed by the
medical-fitness service; and the group chat and alert centre were absorbed by the
notification-mail service, which already owned messages and consumed the events that drive
alerts. That the architecture could grow this much at the leaves while remaining stable at
the trunk is, in the team's assessment, the project's strongest evidence that its founding
design choices were sound.

In short, the system is not a prototype that merely gestures at its intended scope. It is a
runnable, role-aware, secure platform in which the day-to-day operations of a multi-sport
club---squad administration, training, match operations, medical care, recruitment,
commercial negotiation, analytics and fan engagement---are connected through a common
identity and a common data backbone rather than scattered across disconnected tools.

\section{What Was Delivered}
\label{sec:delivered}

The project delivered a complete, runnable multi-sport club management platform whose
principal artefacts are summarised below. The itemisation is deliberately concrete: each
entry names the artefact, its location in the architecture, and the requirement family it
satisfies, so that the claim of completeness can be checked against the requirements
catalogue reproduced in Appendix~A.

\begin{itemize}
  \item \textbf{A microservices back-end.} Seven Spring~Boot services---the six domain
        services user-management, player-management, training-match, medical-fitness,
        reports-analytics and notification-mail, fronted by an API gateway, a configuration
        server and a Eureka service registry---all orchestrated by a single Docker~Compose
        definition alongside Keycloak, PostgreSQL, Kafka, Redis and RabbitMQ. Each domain
        service owns a dedicated database, registers itself for discovery on boot, and
        draws its configuration from the Config Server, so that the entire control plane
        and data plane come up from one reproducible command.

  \item \textbf{Layered identity and access control.} Every request is authenticated
        against the Keycloak \texttt{mscms} realm using the \texttt{keycloakId}
        identifier (the JWT \texttt{sub} claim, never an internal database key), and
        authorisation is enforced consistently at three layers---the identity provider, the
        API gateway and the Next.js route guards---so that the same per-role permission
        rules apply whether a user navigates the interface or calls the APIs directly. The
        front-end permission map is the single source of truth that drives both the
        navigation a user sees and the routes the middleware will admit, eliminating drift
        between what is shown and what is allowed.

  \item \textbf{A real competitions module.} Database-backed competitions modelling
        La~Liga and the new-format UEFA Champions League (a single combined league phase
        followed by a knockout bracket that unlocks only once the league phase concludes),
        with a guided multi-step creation flow, automatic round-robin fixture generation
        for league formats, completed competitions automatically locked as read-only with a
        ``Completed'' badge, and real, finished seasons imported from the football-data.org
        API as seed data. Standings are recomputed from the recorded fixture results rather
        than stored, so the table is always consistent with the matches that produced it.

  \item \textbf{A live-match engine.} A broadcast-style phase machine advancing a match
        through first half, half-time, second half and full time---and, for knockout
        ties, through two periods of extra time and a penalty shoot-out---with the line-up
        locked at creation, phase-bounded and administrator-set stoppage time, per-sport
        substitution limits, in-match event and injury recording, real goal and result
        alerts raised as events occur, and automatic write-back of the final score to the
        source competition fixture. The knockout-versus-points decision is resolved from
        the originating fixture's round, so a draw is correctly permitted for a league-phase
        fixture and correctly disallowed for a knockout tie.

  \item \textbf{A clinical injury-management workflow.} A medical \emph{Treatment Room}
        separating currently-injured from recovered players, with each player advancing
        through a staged journey---injury, diagnosis, treatment, rehabilitation, recovery,
        fitness test and return to fitness---driven by the injury's own status as the single
        source of truth for the current stage. Logging an injury automatically withdraws the
        player from upcoming matches, line-ups and training, and only a passing fitness test
        returns the player to the available pool; the journey is supported by general
        medical reference catalogues for injuries, diagnoses, treatments, rehabilitation,
        recovery programmes and fitness tests.

  \item \textbf{Real-time team messaging.} A group-chat facility built on persistent
        \texttt{ChatGroup} and \texttt{ChatGroupMember} entities, with a seeded general
        group returned to everyone, user-created groups, membership invitations carrying an
        invited-or-member state with explicit accept and decline actions, threaded replies
        and emoji reactions, and client-side de-duplication so each message renders exactly
        once. Delivery uses a native WebSocket channel connected directly to the
        notification service at \texttt{ws://localhost:8085/ws/chat}, deliberately outside
        the gateway's JWT filter, with authenticated polling retained as a resilient
        fallback.

  \item \textbf{An event-driven notification centre.} A top-bar alert bell backed by a
        Kafka consumer that turns domain events---scheduled and completed matches,
        injuries, transfers and cancelled training, together with live-match goal and
        result alerts contributed directly by the live engine---into role-aware,
        priority-tagged notifications. The centre presents alerts and notifications across
        two tabs, carries unread badges, and refreshes on a short polling interval so that
        critical events reach the relevant actor promptly.

  \item \textbf{A polished, role-aware front-end.} A Next.js application providing
        per-role dashboards, a public-facing Club Hub and player profiles with real
        photographs, the competitions and match surfaces described above, aggregated
        analytics and reports, a read-only mode for players and fans, and a settings area
        with real persistence and a Keycloak-verified password-change flow, all presented
        under the \textit{Sportify} brand identity. Auxiliary surfaces such as a head-to-head
        record view and a real FIFA World Cup view, computed entirely from back-end data,
        round out the consumer-facing experience.

  \item \textbf{Project documentation and deployment assets.} Alongside the running code,
        the project delivered an entity-relationship description of the six databases, an
        authentication sequence description, a complete features-and-endpoints reference, a
        Docker quick-start guide, and the prebuilt-image build workflow that allows the
        stack to be assembled reproducibly on a host where in-container Maven cannot reach
        the network. The deployment and run procedure is reproduced in Appendix~E.
\end{itemize}

\section{Lessons Learned}
\label{sec:lessons}

Developing a distributed system of this scope as a student team yielded several
durable lessons, both technical and organisational. They are recorded here not as
platitudes but as concrete observations that changed how the team worked.

\paragraph{Contract drift is the dominant cost of distributed development.}
The single most expensive recurring problem was the drift between the front-end client,
the back-end data-transfer objects and the database schema. A field renamed on one side,
or a status value spelled differently across services, repeatedly produced silent
failures that were costly to trace because no single component was obviously at fault: the
front-end believed it was sending the right shape, the back-end believed it was receiving
the wrong one, and the only evidence was an empty list or a swallowed error. The team
mitigated this by nominating a single source of truth for each cross-cutting concern---for
example a shared permissions module on the front-end and seed-identifier constants on the
back-end---and by treating the REST contract~\cite{restapi} as a deliberate interface
rather than an incidental byproduct of implementation. The broader lesson is that in a
distributed system the interface between components deserves at least as much design care as
the components themselves.

\paragraph{Asynchronous, event-driven design pays for itself as features grow.}
Adopting a Kafka event backbone~\cite{kafka} early meant that later features such as
the notification centre and live-match alerts could be added by consuming events already
being published, rather than by threading new synchronous calls through multiple
services. When the live-match engine needed to announce a goal or a final result, it did
not need to know who cared---it published an event, and the notification service, which
already consumed the domain's event stream, turned it into an alert. This validated the
decision to favour loose coupling and demonstrated, in practice, why event-driven
microservices~\cite{microservices} scale better with feature count than a tightly coupled
design: the cost of adding a consumer is independent of the number of producers.

\paragraph{Centralised identity simplifies everything downstream.}
Delegating authentication to Keycloak~\cite{keycloak} and propagating identity as a
signed JWT~\cite{jwt} removed an entire class of problems---password storage, session
handling and per-service user tables---and made the three-layer RBAC~\cite{rbac} model
tractable. Because every service trusts the same signed token and refers to a person by the
same \texttt{keycloakId}, identity never had to be reconciled across service boundaries.
The lesson was that investing in a single identity provider before building features,
rather than retrofitting one later, repaid the effort many times over; the two-step,
Keycloak-verified password change is a small but telling example of the leverage this
gained, since the same provider that issues tokens also authoritatively re-verifies the
current credential before allowing it to change.

\paragraph{Idempotent, deterministic seeding makes a distributed system safe to operate.}
Guarding every seeding routine so that it runs only against an empty store, and preserving a
deterministic insertion order so that cross-service numeric identifiers remain stable, made
the full stack safe to tear down and redeploy without corrupting data. This proved
indispensable when iterating across seven services with independent databases: the team
could wipe the volumes, bring the stack back up, and rely on a consistent baseline in which
FC~Barcelona always carried the same team identifier in every service that referenced it.

\paragraph{A container-first discipline removes a whole category of integration pain.}
Running every service through Docker~Compose from the earliest iterations~\cite{docker}
eliminated the ``works on my machine'' class of defect and gave every iteration a
reproducible integration environment. The discovery that in-container Maven could not reach
the network on the development machine forced a prebuilt-image workflow---building fat jars
on the host and copying them into minimal runtime images---which, although born of a
constraint, ultimately produced faster, more predictable rebuilds.

\paragraph{Teamwork lessons.}
Working as a multi-person team across loosely coupled services, an agile, iterative
process~\cite{agile} with short cycles and frequent integration was essential: the
team learned to agree interface contracts before parallel work began, to integrate often
rather than at the end, and to keep shared documentation and conventions current so that
a service owned by one member could be safely consumed by another. Clear ownership of
each service, combined with disciplined communication at the integration boundaries,
was the practical key to progressing in parallel without constant blocking. The team also
learned the value of a vertical-slice habit---finishing each iteration with an integrated,
runnable system rather than a drawer of disconnected parts---because it surfaced
integration problems while they were still small and cheap to fix.

\section{Limitations}
\label{sec:limitations}

In the interest of an honest appraisal, the system carries a number of limitations that a
production deployment would need to address. These are stated plainly; none of them
undermines the project's central claim, but each marks a boundary between what was built and
what a fully operational product would require.

\begin{itemize}
  \item \textbf{Single-region, single-host deployment.} The platform runs as a single
        Docker~Compose~\cite{docker} stack on one host. There is no clustering,
        horizontal autoscaling or multi-region redundancy, so the current deployment is
        appropriate for demonstration and evaluation rather than high-availability
        production use. The architecture is designed to permit horizontal scaling---each
        service is stateless behind the gateway and discoverable through Eureka---but that
        capability has not been exercised against a real orchestrator.

  \item \textbf{Manual operational tasks.} Several operational steps remain manual,
        including building service artefacts on the host before image assembly and the
        manual triggering of certain data-synchronisation tasks. There is no continuous
        integration or deployment pipeline, and observability is limited to container
        logs rather than dedicated metrics, dashboards and distributed tracing, which makes
        diagnosing a cross-service problem more laborious than it would be in a mature
        operational setup.

  \item \textbf{Client-side medical reference templates.} The general medical
        catalogue templates that support the injury workflow are partly held on the client
        rather than being fully persisted server-side, which means they are not yet
        centrally managed or shared consistently across sessions and users. The per-player
        clinical records themselves are persisted; it is the reusable template library that
        is incompletely server-backed.

  \item \textbf{Constraints of the external data feed.} Because real fixture data is
        drawn from the free tier of football-data.org~\cite{footballdata}, the system
        is limited to the competitions and seasons that tier exposes; team-scoped endpoints
        are forbidden on the free tier, so the service reads competition-scoped endpoints
        instead. At the project's reference date the available data is post-season and
        therefore consists of finished matches, so the live-match engine is best
        demonstrated on manually created fixtures rather than on genuinely upcoming
        real-world ones.

  \item \textbf{Approximated real-time presence.} Real-time messaging operates over a
        WebSocket channel~\cite{websocket} with a polling fallback, but the system does
        not yet include a dedicated presence service tracking online status, typing
        indicators and delivery and read receipts reliably across all clients. Presence is
        approximated rather than authoritative.

  \item \textbf{Evaluation by inspection rather than automated UI testing.} Validation
        of the user interface was performed largely by manual inspection and by
        screenshots of each role's surfaces (presented in Chapter~\ref{ch:results}),
        rather than by an automated end-to-end test suite, so regression protection on the
        front-end is limited and a future change could silently break a flow that is not
        re-checked by hand.

  \item \textbf{No commercial transaction layer.} The platform models sponsorship and
        transfer workflows but does not implement payment, ticketing or merchandising,
        and therefore does not handle real financial transactions or the compliance,
        reconciliation and fraud concerns that accompany them.

  \item \textbf{Decoupled, snapshot machine learning.} The match-prediction and
        player-rating models are served by standalone FastAPI services reached through a
        proxy rather than the gateway, and they operate on static trained snapshots rather
        than a continuously retrained pipeline, so their predictions do not yet improve
        automatically as fresh match data accumulates.
\end{itemize}

\section{Future Work and Recommendations}
\label{sec:future-work}

The architecture leaves ample room for extension. The following recommendations are
ordered roughly by the value they would add relative to the effort required, and are
grouped into product extensions, which add user-facing capability, and platform
extensions, which raise operational maturity. In every case the recommendation builds on a
mechanism the platform already possesses---the shared identity, the event backbone, the
gateway origin or the reporting service---rather than requiring an architectural departure.

\subsection{Product Extensions}

\begin{itemize}
  \item \textbf{Native mobile application.} A native or cross-platform mobile client
        for players, coaches and fans would reuse the existing gateway and Keycloak
        realm without introducing a new identity surface. Players would receive push
        notifications for line-ups, training schedules and injury updates; coaches would
        log attendance and rate sessions in the field; and an offline-first cache would
        keep the app usable on poor connections at training grounds and stadiums.

  \item \textbf{Payments, ticketing and merchandising.} An e-commerce surface---branded
        kit, season tickets and match-day tickets---wired into the same Keycloak
        identity and fan navigation, owned by a new service with its own database and
        integrated with a payment gateway. Revenue events would flow into the
        reports-analytics service so that commercial performance can be correlated with
        on-pitch results, closing the loop between the sporting and commercial sides of the
        club within a single analytical surface.

  \item \textbf{A true presence and delivery service.} Promoting the current messaging
        layer to a dedicated presence service would add reliable online-status tracking,
        typing indicators and delivery and read receipts over the existing WebSocket
        channel~\cite{websocket}, removing the present reliance on a polling fallback and
        turning approximated presence into authoritative presence.

  \item \textbf{Fan-engagement layer.} Polls, player-of-the-match voting and
        predict-the-score mini-games, all tied to the same fan identity and driven by
        Kafka events~\cite{kafka} already published by the back-end, would deepen
        engagement without adding new write paths to the operational schema. Because the
        live-match engine already emits goal and result events, a predict-the-score game
        could be settled automatically from the event stream.

  \item \textbf{A global clubs and players directory.} Replacing the static reference data
        that currently backs the rival-club directory and the competition team pools with a
        real back-end catalogue would let administrators compose competitions from a
        managed, searchable set of clubs and players, and would enable team and player
        comparison features across the whole directory rather than only the modelled club.

  \item \textbf{Cross-platform search.} A unified search surface spanning players, staff,
        matches, competitions and messages would shorten navigation for power users and is a
        natural addition now that the domains share a common identity and gateway origin.
\end{itemize}

\subsection{Platform Extensions}

\begin{itemize}
  \item \textbf{Deeper AI and machine learning.} The existing match-prediction and
        player-rating models could be promoted from static snapshots to a production
        pipeline that periodically retrains on fresh match data through a versioned
        model registry, with drift alerts delivered through the existing notification
        fabric. Further models---injury-risk estimation from training load, or
        opponent-style classification---would extend the analytical value of the
        platform and would consume data the medical-fitness and training-match services
        already capture.

  \item \textbf{Horizontal scaling and high availability.} Migrating the
        single-host Docker~Compose stack to a container orchestrator with
        per-service horizontal autoscaling would allow the system to absorb fan-traffic
        bursts on match days and to provide the redundancy that production operation
        requires, building naturally on the service decomposition already in
        place~\cite{microservices} and on the stateless, discoverable design of each
        service.

  \item \textbf{Continuous integration, observability and auditing.} Introducing an
        automated build-test-deploy pipeline, centralised metrics and dashboards,
        distributed tracing across services, and a granular audit log of every
        state-changing administrative action would substantially improve operational
        maturity and governance, and would replace the current reliance on raw container
        logs for diagnosis.

  \item \textbf{Automated end-to-end testing.} An automated front-end and integration
        test suite, exercising the key per-role flows against the REST
        contract~\cite{restapi}, would replace much of the present manual inspection
        and provide durable regression protection as the platform evolves; pairing it with
        contract tests at the service boundaries would directly attack the contract-drift
        problem identified among the lessons learned.

  \item \textbf{Hardening for production identity and transport.} Enabling TLS end to
        end, moving refresh-token handling and the WebSocket handshake fully behind the
        gateway, and tightening CORS and cookie flags for an HTTPS deployment would close
        the gap between the demonstration configuration and a production-grade security
        posture, again leveraging Keycloak~\cite{keycloak} as the single identity authority.
\end{itemize}

In summary, MSCMS---delivered under the \textit{Sportify} brand---fulfils its stated
goal of unifying the operations of a modern multi-sport club into a single secure,
cloud-native platform, and its modular, event-driven architecture provides a sound
foundation on which the extensions outlined above can be built incrementally. The project
therefore concludes not as a finished product but as a robust, extensible platform whose
design has already been shown to absorb substantial new behaviour without structural
change---which is precisely the property a long-lived club-management system most needs.
